\documentclass{article}
\usepackage{graphicx} % Required for inserting images
\usepackage{amsmath,amssymb}
\usepackage{pgfplots}
\pgfplotsset{compat=1.18}
\usepackage{tikz}
\usepackage{hyperref}
\title{Three Methods to Compare $e^\pi$ \text{and}\  $\pi^e$}
\author{Sun Shengran}
\date{December 2025}

\pgfplotsset{compat=1.18}
\begin{document}  

\maketitle

\section{Introduction}
It is interesting to combine two essential constants into two numbers and work out
some relationships between them. Thus, in this artice, I will discuss three methods to compare them. The first two methods both relate to logarithm and the last one is more elegant.

\section{First Method: Take Logarithm Both Sides}
\subsection{Lemma: Graph of $y = \frac{\ln(x)}{x},\ x\in \mathbb{R}^{+}$}

\begin{figure}[h]
\centering
\begin{tikzpicture}
  \begin{axis}[
      axis lines=middle,
      xlabel={$x$},
      ylabel={$y$},
      domain=0.1:6,
      samples=200,
      smooth,
      ymin=-2, ymax=1,
      xmin=0, xmax=6,
    ]
     \addplot[blue, thick] {ln(x)/x};
     \addplot[
      red,
      mark=*,
      only marks
     ] coordinates {(e, 1/e)};
     \node[anchor=south] at (axis cs:e, 1/e) {$\left(e,\frac{1}{e}\right)$};
  \end{axis}
\end{tikzpicture}
\end{figure}
\subsection{Main Process of Comparison}

\begin{center}
$e^\pi$\ \text{v.s.}\ $\pi^e$\

\vspace{0.2cm}
$\Leftrightarrow$\  $\pi \ln(e)$\ \text{v.s.}\ $e \ln(\pi)$ \\
\vspace{0.1cm}
\text{(take natural logarithm as logarithm is a strictly increasing function)}

\vspace{0.2cm}
$\Leftrightarrow$\ $\frac{1}{e}$\ \text{v.s.}\ $\frac{\ln(\pi)}{\pi}$
\end{center}
By lemma 2.1, $\frac{1}{e}$ \text{is the maximum of the function}\ $\frac{\ln(x)}{x}$.
Hence, $e^\pi$ $>$ $\pi^e$.


\section{Second Method: A Classic Area Comparison }
Considering the graph of $y = \frac{1}{x},\ x>0$:

\begin{figure}[h]
\centering
\begin{tikzpicture}
  \begin{axis}[
      axis lines=middle,
      xlabel={$x$},
      ylabel={$y$},
      domain=0.5:5,
      samples=200,
      smooth,
      ymin=0, ymax=1,
      xmin=0, xmax=5,
    ]
     \addplot[blue, thick] {1/x};
   \end{axis}
\end{tikzpicture}
\end{figure}

The area of the rectangle with width ($\pi-e$) and height $\frac{1}{e}$
is larger than the area enclosed by the curver $y=\frac{1}{x}$ and x-axis between $x=e$ and $x=\pi$.

Hence,

\begin{center}
$\frac{\pi-e}{e} > \int^{\pi}_{e} \frac{1}{x} dx$

\vspace{0.2cm}
$\Rightarrow$\ $\frac{\pi}{e}-1 > \ln(\pi)-1$

\vspace{0.2cm}
$\Rightarrow$\ $\pi > \ln(\pi^e)$

\vspace{0.2cm}
$\Rightarrow$\ $e^\pi > \pi^e$.

\end{center}


\section{Third Method: Applying $e^x \ge x+1$}
\subsection{Lemma: $e^x \ge x+1,$ \ attaining equal sign iff $x=0$}
This can be simply inspected from the graphs of $y=e^x$\ and $y=x+1$ on the same axes.

\subsection{Main Process of Comparison}
\begin{center}
$e^\pi = e^{e(\frac{\pi}{e})} = (e^{(\frac{\pi}{e}-1+1)})^e = (e^{(\frac{\pi}{e}-1)})^e e^e$
$\ge (\frac{\pi}{e})^e e^e = \pi^e$

\vspace{0.5cm}
as ${\frac{\pi}{e}-1} \neq 0$,\ $e^\pi > \pi^e$.
\end{center}


\section{Summary}
The first method involves taking logarithm which is the most common approach to tackle exponents. Using the graph of $y = \frac{\ln(x)}{x}$ which we are familiar with is actually a kind of symmetry exploitation.  

In comparison,the second method seems to be not that natural and not easy notice. Although this approach, relying on the area under the curve $y = \frac{1}{x}$ is quite common, it is not simple to come up with it immediately.

The third method is more competition-styled. The application of the inequality $e^x \ge x+1$ is not uncommon in mathematics competition. Even so, the key point of this approach is the algebraic manipulation that is incredibly grateful and skillful.

\section{Reference}
1. Second Method Source: \ \href{https://pmt.physicsandmathstutor.com/download/Maths/A-level/Papers/OCR-MEI/Paper-3/IN/Nov%202020%20IN.pdf}{A level Mathematics B (MEI), October 2021, OCR}
\end{document}
